\chapter{Conclusion}
\label{chap:conclusion}

\section{Summary of Results}
\label{sec:conc_summary}

We asked a single practical question: how many FEM solver calls can be replaced
by PINN predictions in a transient water quenching simulation without losing
useful accuracy?
To answer it, we built the MSWP framework, selected three NAS architectures,
and evaluated the framework across ten unique benchmark geometries.
All domains used the physical parameters of Mortensen et al.~\cite{mortensen2026}.

The answer was yes, and at a higher skip factor than we initially expected.
We found that Bayesian/TPE stays below $5\,^\circ\mathrm{C}$ through $k=3$
on all three 2D canonical domains.
At $k=5$, Rectangle and L-Shape remain below $6\,^\circ\mathrm{C}$.
Circle rises to about $10\,^\circ\mathrm{C}$ at $k=5$, showing that the maximum
useful skip factor is geometry-dependent.
In 3D, Bayesian/TPE reached $k=5$ on all four canonical geometries with errors
below $14\,^\circ\mathrm{C}$, replacing 16 of 20 FEM solver calls.
On the Thermal Fin benchmark, Bayesian/TPE with Fourier embedding stayed below
the $15\,^\circ\mathrm{C}$ threshold at all five skip factors.
We found that NSGA-II and NSGA-III are reliable through $k=3$--$4$ on most
geometries.

We also evaluated the adaptive-$k$ strategy on four 2D domains.
All four architectures achieved 83--85\% FEM saving with MAE below
$2.5\,^\circ\mathrm{C}$.
We concluded that adaptive-$k$ is the most appropriate strategy for geometrically
simple domains where the temperature gradient is steepest in the early phase
of quenching.

We found that FEM anchoring cannot be removed without a large accuracy cost.
In PINN-only mode, mean-window errors were between three and eighteen times
higher than in the anchored hybrid.
On the Thermal Fin the best PINN-only result was $59.20\,^\circ\mathrm{C}$,
which is four times the acceptance threshold
(Section~\ref{subsec:tf_pinn_only}).

We observed that the maximum acceptable skip factor is geometry-dependent.
There is no single $k$ that is optimal across all domains.
A short calibration run on the target geometry is required before deployment.
We also found that mean-window MAE is the appropriate primary metric.
Final-window MAE is unreliable because the temperature field approaches
uniformity at $t=30\,\mathrm{s}$ regardless of intermediate prediction quality.

Every experiment in this thesis was designed with the Mortensen et al.\
industrial quenching simulation~\cite{mortensen2026} as the target deployment
context.
We used the physical parameters and engineering tolerance levels from that work
directly.
We concluded that the MSWP framework is ready for the next phase: applying the
validated method to the real Mortensen mesh with the temperature-dependent
boiling curve.
Section~\ref{sec:conc_future} describes the steps that separate the current
implementation from that goal.

\section{Limitations}
\label{sec:conc_limitations}

We used constant material properties throughout.
Real A356 aluminium has temperature-dependent thermal conductivity and undergoes
latent heat release during solid-state phase changes.
This simplification is standard in benchmark PINN studies but overestimates
cooling uniformity in the early high-temperature phase.
We also used a constant convection coefficient.
Industrial water quenching transitions through several boiling regimes as the
surface cools below the Leidenfrost point.
Our constant-$h$ model does not capture this transition.

All reference solutions came from our own Crank--Nicolson solver.
We did not validate against a commercial code such as Abaqus.
Mesh independence was also not verified.
All grids were fixed structured meshes.
Whether reported MAE values change on finer grids remains untested.

We ran ten seeds per configuration.
Variance remained small at low skip factors and widened slightly at $k=5$ on
the most demanding geometries.
The NAS search used 150 trials for Bayesian/TPE and 15 generations for
NSGA-II/III.
Larger budgets could produce different architectures, especially on geometries
that differ substantially from the 2D rectangle used as the NAS validation domain.

We note that linear interpolation between FEM snapshots achieves lower $L_2$
error than the trained PINN on smooth canonical domains.
This is physically expected since exponential cooling is nearly linear in time
between closely spaced anchors.
The PINN provides value through physics-consistent predictions, fast inference
after training, and adaptability to complex spatial structure.
Identifying the geometry complexity threshold at which PINN surpasses
interpolation is an open question for future work.

\section{Future Work}
\label{sec:conc_future}

We ordered the following steps as a deployment roadmap.
Each brings the MSWP framework one stage closer to the full Mortensen et al.\
industrial simulation~\cite{mortensen2026}.
The framework was designed so that each extension can be added independently
without restructuring the core prediction pipeline.

The most immediate step is applying Bayesian/TPE directly to the Mortensen
industrial mesh and FEM snapshots.
We designed the codebase to support this without code changes.
The data loader accepts any NPZ grid, the domain sampler adapts automatically
to 3D geometries, and the boiling-curve module already has a placeholder for
piecewise $h(T)$.
This step would produce the first direct comparison between PINN-predicted
temperatures and the validated distortion predictions from~\cite{mortensen2026}.

The Mortensen simulation uses a full five-regime temperature-dependent boiling
curve.
Incorporating this requires three changes: the FEM solver must adopt a nonlinear
solve at each step, the PINN boundary residual must use the same $h(T)$ function,
and the boiling-curve coefficients must be sourced from~\cite{mortensen2026}.
A prototype of the boundary residual is already included in the codebase.
Once in place the core prediction loop is unchanged.

The adaptive-$k$ controller we tested used a single promotion threshold for all
architectures.
We found that this causes a threshold mismatch because the three architectures
operate at different baseline error levels.
A deployment-ready controller should initialise thresholds from the known
fixed-$k$ optimum of each architecture separately.
This would allow higher FEM savings than the static-threshold version we evaluated.

We used uniform collocation sampling inside each domain.
A residual-based sampler would concentrate training effort where the physics is
hardest to satisfy.
This would generalise to complex geometries without manual tuning and is
especially relevant for the irregular surfaces of the Mortensen mesh.

We generated all FEM reference solutions with our own solver.
Validating a representative subset of benchmarks against StaMiSim or Abaqus
would confirm that the reported PINN errors reflect true physics accuracy and
are not artefacts of the in-house solver.

\section{Closing Remarks}
\label{sec:conc_closing}

Mortensen et al.~\cite{mortensen2026} solve the transient heat equation for
water quenching of A356 aluminium subframes using full FEM at every time step.
We proposed solving the same physical problem with NAS-guided PINNs as the
primary prediction engine and retaining FEM only as a periodic correction.
We answered the central question across ten unique benchmark geometries.

The answer is geometry-dependent but consistently positive.
We found that 65--85\% of FEM calls can be eliminated while keeping mean-window
temperature errors within practical bounds.
At $k=3$ this holds for all eight canonical geometries.
At $k=4$ and $k=5$ it holds for the majority, with Bayesian/TPE as the most
robust choice.

At the same time, we found that the method is not ready to run without FEM calls.
The PINN-only results showed clearly that the network needs periodic corrections,
particularly in the early part of the simulation when the cooling rate is highest.
Reducing the dependence on FEM anchoring while maintaining accuracy is the
central open question for this line of research.

The future work steps in Section~\ref{sec:conc_future} form a concrete roadmap
from the current benchmark-validated implementation to a full deployment on the
Mortensen et al.\ industrial quenching simulation.
The results presented here are not a standalone benchmark study.
They are the validated foundation for a PINN-based computational accelerator
for industrial aluminium quenching.
